\subsubsection{Программное обеспечение}\label{subsec:ch1/sec2/sub4/sub3}

Программное обеспечение является основой для автоматизации НТТС. Оно позволяет выполнять широкий спектр задач: от управления движением до взаимодействия с внешней средой. Эффективность работы НТТС напрямую зависит от качества используемого ПО, которое делится на несколько категорий (рис.\:\cref{fig:rms_drawio}).

\begin{figure}[ht]
	\centerfloat{
		\includegraphics[width=14cm]{drawio/rms.drawio.png}
	}
	\caption{Программная система управления НТТС}\label{fig:rms_drawio}
\end{figure}

ПО для управления транспортными средствами \cite{LinuxKetov, LinuxWard}. Данная категория включает системы, которые обеспечивают непосредственное управление транспортным средством, включая управление движением, манипуляторами и исполнительными механизмами.

Примеры:

Системы автоматического вождения (например, автопилоты для карьерных самосвалов).
Контроллеры манипуляторов для погрузочно-разгрузочных операций.
Эти системы часто применяются на уровне отдельных транспортных средств и имеют ограниченную функциональность по взаимодействию с другими элементами системы.
ПО для управления процессами
Эти системы отвечают за координацию и управление задачами на уровне всего производственного процесса.
Направления:

Системы планирования маршрутов для транспортных средств в логистических цепочках.
Системы диспетчеризации, которые распределяют задачи между машинами и контролируют их выполнение.
Примером является система управления автопарком (Fleet Management System, FMS), применяемая в логистике и горнодобывающей промышленности.
ПО для анализа и прогнозирования
Данные системы используют машинное обучение и методы анализа данных для прогнозирования и оптимизации работы НТТС.
Примеры:

Прогнозирование состояния транспортных средств для предотвращения поломок (predictive maintenance).
Оптимизация потребления топлива на основе анализа маршрутов и условий работы.
Такие системы часто интегрируются с облачными платформами для доступа к большому объёму данных.
ПО для взаимодействия в многоагентных системах
Данная категория фокусируется на создании координированных действий между несколькими участниками процесса.
Направления:

Системы, обеспечивающие координацию автономных транспортных средств в рамках одного объекта (например, логистического склада).
ПО для взаимодействия автономных транспортных средств с внешней инфраструктурой, включая системы связи V2X (vehicle-to-everything).
Перспективные направления разработки ПО
Интеграция искусственного интеллекта (ИИ)

Использование ИИ для динамического перераспределения задач между участниками.
Реализация систем глубокого обучения для адаптации к меняющимся условиям работы.
Облачные платформы для управления процессами

Централизованное управление данными и задачами через облачные системы.
Интеграция различных типов ПО для создания единого цифрового пространства.
Системы коллективного управления
Перспективными являются системы, которые обеспечивают совместное выполнение задач группой автономных транспортных средств. Например, автономные конвои для логистики или синхронная работа транспортных средств на карьерах.

Универсальные стандарты и открытые платформы

Создание стандартизированных решений, которые позволяют разным системам взаимодействовать друг с другом.
Открытые программные интерфейсы (API) для упрощения интеграции ПО.